\subsection{Criteria for measurability}

When F is a topological vector space, every step function over the measurable sets (§4, No.~9, Def. 4), with values in F, is obviously measurable (No.~3, Cor. 3 of Th. 1); such a function f takes on only a finite number of values, and for every \(y \in F\) , \(f(y)\) is measurable. More generally, let F be any topological space, f a mapping of X into F taking on only a finite number of distinct values \(a_i\) ( \(1 \leqslant i \leqslant m\) ); if the sets \(A_i = f(a_i)\) are measurable, then the function f is measurable. For, for every compact set K and for each of the sets \(A_i \cap K\) , there exist a negligible set \(N_i \subset A_i \cap K\) and a partition of \(A_i \cap K \cap C N_i\) formed by a sequence ( \(K_{in}\) ) of compact sets; since K is the union of the sets \(A_i \cap K\) and the restriction of f to each of the \(K_{in}\) is constant, therefore continuous, f is measurable. By an abuse of language we shall say that a mapping f of X into F is a measurable step function if it takes on only a finite number of values and if, for every \(y \in F\) , \(f(y)\) is measurable.

THEOREM 3. --- For a mapping f of X into a metrizable space F to be measurable, it is necessary and sufficient that, for every compact set \(K \subset X\) , there exist a sequence \((g_{n})\) of measurable step functions, with values in F, such that \(g_{n}(x)\) tends to \(f(x)\) for almost every \(x \in K\) .

The condition is sufficient by Egoroff's theorem and the principle of localization. Let us show that it is necessary: by hypothesis, there exist a negligible set \(N \subset K\) and a partition \((K_{m})\) of K - N formed of compact sets such that the restriction of f to each \(K_{m}\) is continuous. To define the sequence \((g_{n})\) , it suffices to proceed in the following manner: let d be a metric compatible with the topology of F; for each \(K_{i}\) with index \(i \leqslant n\) , there exists a finite partition of \(K_{i}\) into integrable sets \(A_{ij}\) ( \(1 \leqslant j \leqslant q_{i}\) )

sufficiently small that the oscillation of \(f\) on each of the \(A_{ij}\) is \(\leqslant 1/n\) (§4, No.~9, Lemma); take \(g_n\) to be constant on each \(A_{ij}\) , equal to one of the values of \(f\) in this set (for \(1 \leqslant i \leqslant n\) , \(1 \leqslant j \leqslant q_i\) ), and equal to a fixed element \(a \in F\) for every point of \(X\) that does not belong to any of the \(A_{ij}\) . It is clear that the sequence \((g_n(x))\) converges to \(f(x)\) at every point of \(K\) not belonging to \(N\) .

COROLLARY 1. --- Let f be a measurable mapping of X into a Banach space F; for every compact set K ⊂ X, there exists a sequence (gₙ) of measurable step functions, with supports contained in K, such that \textbar gₙ(x)\textbar{} ≤ \textbar f(x)\textbar{} for all x ∈ X and such that gₙ(x) tends to f(x) for almost every x ∈ K.

With notations as in the proof of Th. 3, and denoting by \(\mathbf{a}_{ij}\) one of the values of \(\mathbf{f}\) in \(\mathrm{A}_{ij}\) , it suffices to take, as the value of \(\mathbf{g}_n\) in \(\mathrm{A}_{ij}\) , the point 0 if \(|\mathbf{a}_{ij}| \leqslant 1/n\) and the point \(\mathbf{a}_{ij}(1 - 1/(n|\mathbf{a}_{ij}|))\) otherwise; finally, set \(\mathbf{g}_n(x) = 0\) on the complement of the union of the \(\mathrm{A}_{ij}\) ( \(1 \leqslant i \leqslant n\) , \(1 \leqslant j \leqslant q_i\) ).

COROLLARY 2. --- Let X be a locally compact space that is countable at infinity. If f is a measurable mapping of X into a metrizable space F, then there exists a sequence \((g_{n})\) of measurable step functions, with values in F, such that \(g_{n}(x)\) tends to \(f(x)\) for almost every \(x \in X\) .

If X is the union of an increasing sequence \((\mathbf{A}_{n})\) of compact sets, the sets \(A_{n}-A_{n-1}\) that are nonempty form a partition of X into integrable sets; consequently, there exist a negligible set \(N\subset X\) and a partition of CN formed by a sequence \((\mathrm{K}_{n})\) of compact sets such that the restriction of f to each \(K_{n}\) is continuous; the proof can then be concluded as in Th. 3 without modification.

PROPOSITION 7. --- Let f be a measurable mapping of X into a topological space F; the inverse image under f of every closed (resp. open) set in F is measurable.

It suffices to carry out the proof for the inverse image \(f(A)\) of a closed set A in F. Let K be a compact subset of X; there exist a negligible set \(N \subset K\) and a partition \((K_n)\) of K - N formed of compact sets such that the restriction of f to each \(K_n\) is continuous. The intersection \(K_n \cap f(A)\) is thus the inverse image of the closed set A under the restriction of f to \(K_n\) ; it is therefore a closed set in \(K_n\) , hence is compact. The set \(K \cap f(A)\) is therefore the union of the negligible set \(N \cap f(A)\) and the compact sets \(K_n \cap f(A)\) , which proves that \(f(A)\) is measurable.

THEOREM 4. --- Let F be a metrizable space and let d be a metric on F compatible with the topology. For a mapping f of X into F to be measurable, it is necessary and sufficient that it satisfy the following two conditions:

\begin{enumerate}
\def\labelenumi{\alph{enumi})}
\item
  the inverse image under \(f\) of every closed ball of \(F\) is measurable;
\item
  for every compact set \(\mathrm{K} \subset \mathrm{X}\) , there exists a countable subset \(\mathrm{H}\) of \(\mathrm{F}\) such that \(f(x) \in \overline{\mathrm{H}}\) for almost every \(x \in \mathrm{K}\) .
\end{enumerate}

The condition a) is necessary by Prop. 7; on the other hand, in the notations of Th. 3, condition b) is satisfied by taking H to be the countable set formed by the values of all of the functions \(g_{n}\) .

Let us now show that the conditions a) and b) are sufficient. Let K be any compact subset of X; there exists a negligible subset N of K such that \(f(\mathrm{K}-\mathrm{N})\) is contained in the closure of a countable set of points of F, which we arrange in a sequence \((a_{n})\) . Let \(A_{n,p}\) be the set of \(x \in K-N\) such that \(d(f(x), a_{n}) \leqslant 1/p\) . It follows from condition a) that \(A_{n,p}\) is measurable. For fixed p, define a sequence of sets \(B_{n,p} \subset K-N\) recursively by setting

\[
\mathrm{B} _ {1, p} = \mathrm{A} _ {1, p} \quad \text {and} \quad \mathrm{B} _ {n + 1, p} = \mathrm{A} _ {n + 1, p} \cap \mathbb {C} \Big (\bigcup_ {k \leqslant n} \mathrm{A} _ {k, p} \Big);
\]

the sets \(B_{n,p}\) are measurable, and those that are nonempty form a partition of K-N. Let \(g_{m,p}\) be the function equal to \(a_{i}\) on the set \(B_{i,p}\) for \(1 \leqslant i \leqslant m\) and equal to a constant \(b \in F\) on the complement of the union of these sets; \(g_{m,p}\) is a measurable step function; as m tends to infinity, \(g_{m,p}\) converges pointwise to the function \(f_{p}\) equal to \(a_{n}\) on \(B_{n,p}\) ( \(n \geqslant 1\) ) and to b on N∪CK, therefore (Th. 2) \(f_{p}\) is measurable. As p tends to infinity, \(f_{p}(x)\) tends to \(f(x)\) for every \(x \in K-N\) , and to b for \(x \in N \cup CK\) ; the limit of the \(f_{p}\) is therefore measurable, and the principle of localization proves that f itself is measurable.

Remarks. --- 1) Condition a) alone is not sufficient for \(f\) to be measurable (Exer. 7).

For, if this is so, for every \(b \in \overline{R}\) the set of x such that \(f(x) \geqslant b\) is measurable, being the intersection of the sets (which form a countable family) of the x such that \(f(x) \geqslant a\) , as a runs over the set of points of D that are \(\leqslant b\) . The set of x such that \(f(x) < b\) is measurable, being the complement of a measurable set. Next, if b is finite then the set of x such that \(f(x) \leqslant b\) is measurable, being the intersection of the sets of the x such that \(f(x) < b + 1/n\) ; and \(\overline{f}(-\infty)\) is measurable, being the intersection of the sets of the x such that \(f(x) < n\) , as n runs over Z. Finally, the inverse image under f of every closed interval of \(\overline{R}\) is measurable, being the intersection of two measurable sets, and Th. 4 may be applied.

One could show similarly that it suffices that for every \(a \in D\) , the set of \(x\) such that \(f(x) > a\) be measurable.

COROLLARY. --- Every lower (resp. upper) semi-continuous function is measurable.

For, if \(f\) is lower semi-continuous, then the set of \(x \in \mathbf{X}\) such that \(f(x) \leqslant a\) is closed for every \(a \in \overline{\mathbf{R}}\) .

When F is metrizable and compact, Prop. 7 makes it possible to sharpen Th. 3 as follows:

PROPOSITION 9. --- If F is a metrizable compact space, then every measurable mapping f of X into F is the uniform limit (on all of X) of a sequence of measurable step functions.

Let \(d\) be a metric compatible with the topology of \(\mathbf{F}\) . For every positive integer \(n\) there exists a finite number of points \(a_{k} \in \mathbf{F}\) such that the closed balls \(\mathbf{B}_{k}\) with center \(a_{k}\) and radius \(1/n\) form a covering of \(\mathbf{F}\) ; the sets \(\mathbf{A}_{k} = f(\mathbf{B}_{k})^{-1}\) are therefore measurable (Prop. 7) and form a covering of \(\mathbf{X}\) . Consequently (§4, No.~9, Lemma) there exists a partition \((\mathbf{C}_{i})\) of \(\mathbf{X}\) into a finite number of measurable sets such that each \(\mathbf{A}_{k}\) is the union of certain of the \(\mathbf{C}_{i}\) . Let \(c_{i}\) be a point of \(\mathbf{C}_{i}\) and let \(g_{n}\) be the measurable step function equal to \(f(c_{i})\) on \(\mathbf{C}_{i}\) (for each index \(i\) ). It is clear that \(d(f(x), g_{n}(x)) \leqslant 2/n\) for all \(x \in \mathbf{X}\) .

PROPOSITION 10. --- Let F be a separable Banach space, F' its dual, and (a'\_n) a weakly dense sequence in the unit ball of F' (TVS, III, §3, No.~4, Cor. 2 of Prop. 6).¹ For a mapping f of X into F to be measurable, it is necessary and sufficient that for every n the scalar function x ↦ ⟨f(x), a'\_n⟩ be measurable.

The condition being obviously necessary (No.~3, Th. 1), let us prove that it is sufficient; it suffices to verify condition a) of Th. 4 and, to this end, it will suffice by the principle of localization to prove that, for every compact subset K of X and every closed ball \(B \subset F\) , with center a and radius r, the set \(A = K \cap \mathbf{f}^{-1}(B)\) is measurable; now, for every \(z \in F\) ,

\[
| \mathbf {z} | = \sup _ {n} | \langle \mathbf {z}, \mathbf {a} _ {n} ^ {\prime} \rangle | / | \mathbf {a} _ {n} ^ {\prime} |;
\]

A is thus the intersection of \(K\) and the sets defined by

\[
\left| \left\langle \mathbf {f} (x), \mathbf {a} _ {n} ^ {\prime} \right\rangle - \left\langle \mathbf {a}, \mathbf {a} _ {n} ^ {\prime} \right\rangle \right| \leqslant r \left| \mathbf {a} _ {n} ^ {\prime} \right|;
\]

as these sets are measurable by hypothesis, A is measurable.

COROLLARY 1. --- Let F be a Banach space. In order that a mapping f of X into F be measurable, it is necessary and sufficient that it satisfy the following two conditions:

\begin{enumerate}
\def\labelenumi{\alph{enumi})}
\tightlist
\item
  for every \(\mathbf{a}' \in \mathrm{F}'\) , the scalar function \(x \mapsto \langle \mathbf{f}(x), \mathbf{a}' \rangle\) is measurable;
\item
  for every compact set \(\mathrm{K} \subset \mathrm{X}\) , there exists a countable subset \(\mathrm{H}\) of \(\mathrm{F}\) such that \(\mathbf{f}(x) \in \overline{\mathrm{H}}\) for almost every \(x \in \mathrm{K}\) .
\end{enumerate}

The necessity of the conditions follows from No.~3, Th. 1, and Th. 4. To prove that the conditions are sufficient, it again suffices to verify condition a) of Th. 4. With notations as in the proof of Prop. 10, we may (on account of b)) suppose, after modifying f if necessary on a negligible set, that \(\mathbf{f}(\mathrm{K}) \subset \overline{\mathrm{H}}\) , where H is a countable subset of F. If V is the closed linear subspace of F generated by the set \(H \cup \{a\}\) , then V is a separable Banach space and every continuous linear form on V is the restriction of a form \(a' \in F'\) ; the same reasoning as in Prop. 10 then shows that \(K \cap \overline{\mathbf{f}}(B)\) is measurable.

COROLLARY 2. --- Let F be a locally convex space that is metrizable and separable, and let \(F'\) be its dual. For a mapping f of X into F to be measurable, it is necessary and sufficient that for every \(a' \in F'\) , the scalar function \(x \mapsto \langle \mathbf{f}(x), \mathbf{a}' \rangle\) be measurable.

We may regard F as a subspace of a countable product \(\prod_{n}E_{n}\) of Banach spaces (TVS, II, §4, No.~3), and we can suppose that \(\operatorname{pr}_{n}(\mathbf{F})\) is dense in \(E_{n}\) , which is therefore separable. For every n, the mapping \(pr_{n}\circ f\) is then measurable by Prop. 10, therefore f is measurable by No.~3, Th. 1.

PROPOSITION 11. --- Let F be a locally convex space that is the direct limit of a sequence of separable metrizable locally convex spaces \(F_{n}\) , F being the union of the \(F_{n}\) . Let \(F'\) be the dual of F, equipped with the weak topology \(\sigma(\mathrm{F}',\mathrm{F})\) . For a mapping f of X into \(F'\) to be measurable, it is necessary and sufficient that, for every \(a \in F\) , the scalar function \(x \mapsto \langle \mathbf{a}, \mathbf{f}(x) \rangle\) be measurable.

The condition is necessary by No.~3, Th. 1; let us prove that it is sufficient. Suppose first that F is metrizable and separable, and let D be a countable dense set in F. Let \((\mathrm{V}_{n})\) be a decreasing fundamental sequence of balanced convex open neighborhoods of 0 in F; the polar sets \(V_{n}^{\circ}\) are equicontinuous and their union is all of \(F'\) . Let \(X_{n} = \overline{\mathbf{f}}(V_{n}^{\circ})\) ; the sequence \((\mathrm{X}_{n})\) is increasing and \(X = \bigcup_{n} X_{n}\) ; let us show that each \(X_{n}\) is \(\mu\) -measurable. Indeed, \(D \cap V_{n}\) is dense in \(V_{n}\) ; for every \(y \in D \cap V_{n}\) let \(S_{y}\) be the set of \(x \in X\) such that \(|\langle y, f(x) \rangle| \leqslant 1\) ; the hypothesis implies that each of the \(S_{y}\) is measurable, and \(X_{n}\) is the intersection of the countable family of \(S_{y}\) for \(y \in D \cap V_{n}\) . This being so, for every compact subset K of X and every \(\varepsilon > 0\) , there exists an integer n such that \(|\mu| (K - (K \cap X_{n})) \leqslant \varepsilon/4\) , and then a compact subset \(K_{1}\) of \(K \cap X_{n}\) such that \(|\mu| ((K \cap X_{n}) - K_{1}) \leqslant \varepsilon/4\) ; finally, there exists a compact subset \(K_{2}\) of \(K_{1}\) such that \(|\mu| (K_{1} - K_{2}) \leqslant \varepsilon/2\) and such that the restrictions to \(K_{2}\) of all of the functions \(\langle y, f \rangle\) , where \(y \in D\) , are continuous (No.~1, Prop. 2). Since the set \(f(K_{2}) \subset f(X_{n}) \subset V_{n}^{\circ}\) is equicontinuous, the topology induced by \(\sigma(F', F)\) on \(f(K_{2})\) is identical to the topology of pointwise convergence in D (GT, X, §2, No.~4, Th. 1); the restriction of f to \(K_{2}\) is therefore continuous, whence our assertion in this first case.

Let us pass to the general case. If \(\mathbf{z}'\) is a continuous linear form on \(\mathbf{F}\) , its restriction \(\mathbf{z}_n'\) to \(\mathbf{F}_n\) is continuous; since \(\mathbf{F} = \bigcup_{n} \mathbf{F}_n\) , the dual \(\mathbf{F}'\) of \(\mathbf{F}\) may be identified (algebraically) with a linear subspace of the product \(\prod_{n} \mathbf{F}_n'\) , and then \(\operatorname{pr}_n(\mathbf{z}') = \mathbf{z}_n'\) . Moreover, since each finite subset of \(\mathbf{F}\) is contained in one of the \(\mathbf{F}_n\) , the topology \(\sigma(\mathbf{F}', \mathbf{F})\) is none other than the topology induced by the product topology of the topologies \(\sigma(\mathbf{F}_n', \mathbf{F}_n)\) . This being so, if \(\langle \mathbf{a}, \mathbf{f} \rangle\) is measurable for every \(\mathbf{a} \in \mathbf{F}\) , then so is \(\langle \mathbf{a}_n, \operatorname{pr}_n \circ \mathbf{f} \rangle\) for every \(n\) and every \(\mathbf{a}_n \in \mathbf{F}_n\) , since \(\langle \mathbf{a}_n, \operatorname{pr}_n \circ \mathbf{f} \rangle = \langle \mathbf{a}_n, \mathbf{f} \rangle\) ; the first part of the proof then shows that \(\operatorname{pr}_n \circ \mathbf{f}\) is measurable for every \(n\) , therefore so is \(\mathbf{f}\) (No.~3, Th. 1).

